这个长正合列就是 $\mathbb V(\underline E)$ 上 $\underline O_{e(S)}$ 的 Koszul 解消在 $\mathbb V(\underline E)^*$ 上的限制.

经仿射态射 $q$ 将长正合列 (1.1.5) 直像到 $\mathbb P(\underline E)$. 由 (1.1.1) 和公式 $q_*{a'}^*=q_*q^*p^*$, 得到长正合列

\begin{equation}\tag{1.1.6}
\cdots\longrightarrow
\bigoplus_{n\in\mathbb Z}p^*\bigwedge^i\underline E(n)
\xrightarrow{\ \iota_X\ }
\bigoplus_{n\in\mathbb Z}p^*\bigwedge^{i-1}\underline E(n)
\longrightarrow\cdots .
\end{equation}

它分解为短正合列, 即 (1.1.3) 的像,

\begin{equation}\tag{1.1.7}
0\longrightarrow
\bigoplus_{n\in\mathbb Z}\Omega^i_{\mathbb P(\underline E)/S}(n)
\longrightarrow
\bigoplus_{n\in\mathbb Z}p^*\bigwedge^i\underline E(n)
\xrightarrow{\ \iota_X\ }
\bigoplus_{n\in\mathbb Z}\Omega^{i-1}_{\mathbb P(\underline E)/S}(n)
\longrightarrow0.
\end{equation}

正合列 (1.1.7) 是下列正合列的直和

\begin{equation}\tag{1.1.8}
0\longrightarrow\Omega^i_{\mathbb P(\underline E)/S}(n)
\longrightarrow p^*\bigwedge^i\underline E(n-i)
\xrightarrow{\ \iota_X\ }
\Omega^{i-1}_{\mathbb P(\underline E)/S}(n)
\longrightarrow0.
\end{equation}

回顾 (EGA III, 2.1.15, 或 FAC), 层 $R^ip_*\underline O(k)$ 局部自由, 其构造与任意基变换相容, 且

\begin{equation}\tag{1.1.9}
R^ip_*\underline O(k)=0
\quad\text{若 }i\ne0,r,
\quad\text{或若 }i=0, k<0,
\quad\text{或若 }i=r, k>-r-1,
\end{equation}

因而 $R^ip_*p^*\bigwedge^j\underline E(k)$ 在相同条件下为零.

正合列 (1.1.8) 给出 $\Omega^i_{\mathbb P(\underline E)/S}(n)$ 的解消

\begin{equation}\tag{1.1.10}
0\longrightarrow\Omega^i_{\mathbb P(\underline E)/S}(n)
\xrightarrow{\ E\ }p^*\bigwedge^i\underline E(n-i)
\xrightarrow{\ \iota_X\ }\cdots\longrightarrow
p^*\bigwedge^0\underline E(n)\longrightarrow0
\end{equation}

由 (1.1.9), 这个解消定义的谱序列由于次数的原因而退化 ($E_1=E_\infty$); 对 $0<j<r$, 它给出

\begin{equation}\tag{1.1.11}
R^jp_*\Omega^i_{\mathbb P(\underline E)/S}(n)
=\mathcal H\left(\cdots\xrightarrow{\ \iota_X\ }
p_*\bigwedge^{i-j}\underline E(n-i+j)
\xrightarrow{\ \iota_X\ }\cdots\right).
\end{equation}

由于 $r\geq1$, 由深度论证可知

\begin{equation}\tag{1.1.12}
\bigoplus_{n\in\mathbb Z}p_*p^*\bigwedge^i\underline E(n)
=p_*q_*q^*p^*\bigwedge^i\underline E
={a'}_*{a'}^*\bigwedge^i\underline E
=a_*a^*\bigwedge^i\underline E.
\end{equation}

%%%%%%%%%% source scan idx 50 | folio 43 %%%%%%%%%%

于是 $p_*p^*\bigwedge^i\underline E(n)$ 是 $a_*a^*\bigwedge^i\underline E$ 中权为 $n+i$ 的齐次部分, 即 $\operatorname{Sym}^n(\underline E)\otimes\bigwedge^i\underline E$. 考虑到 (1.1.5) 来自 Koszul 解消, (1.1.11) 表明: 对 $0<j<r$, 层 $R^jp_*\Omega^i_{\mathbb P(\underline E)/S}(n)$ 只有在 $i-j=0$ 且 $n-i+j=0$, 即 $n=0$ 且 $i=j$ 时才可能非零. 此外, 在此情形下, $R^ip_*\Omega^i_{\mathbb P(\underline E)/S}$ 同构于 $\underline O_S$.

对于 $j=0$, (1.1.10) 还给出

\begin{equation}\tag{1.1.13}
p_*\Omega^i_{\mathbb P(\underline E)/S}(n)
=\ker\!\left(\iota_X:
p_*\bigwedge^i\underline E(n-i)\longrightarrow
p_*\bigwedge^{i-1}\underline E(n-i+1)\right),
\end{equation}

同一论证表明, 当 $n-i\leq0$ 时 $\iota_X$ 为单射.

在 (1.1.8) 中取 $n=0$ 和 $i=r+1$, 得到

\begin{equation*}
\Omega^r_{\mathbb P(\underline E)/S}\simeq
p^*\bigwedge^{r+1}\underline E(-r-1).
\end{equation*}

当 $S$ 为域的谱时, 1.1(ii) 和 (iii) 中的消失断言在 $j\ne r$ 时由上述论证得到, 在 $j=r$ 时则通过对偶性 ([2], V, 11.2, 第~211 页) 从上述论证得到. 特别地, 对任意 $j$ 和 $n$, 至多有一个 $i$ 使 $R^ip_*\Omega^j_{\mathbb P(\underline E)/S}(n)\ne0$. 由于 $\Omega^j_{\mathbb P(\underline E)/S}$ 在 $S$ 上平坦, 这就证明了 (i) (最困难的是在 EGA III 中找到出处; 我们建议使用 6.10.5, 7.4.1 和 7.5.5, 然后取极限).

已知对 $0\leq i\leq r$, $R^ip_*\Omega^i_{\mathbb P(\underline E)/S}$ 为可逆层, 只须证明当 $S$ 为域的谱时 $\eta^i\ne0$. 实际上, $\eta^r$ 是一个有理点 (即 $r$ 个超平面的交) 的上同调类, 因而是 $H^r(\mathbb P^r,\Omega^r)$ 的一组基 (FGA, 第 149 讲, 第~4 节).

1.2. $\Omega^i(n)$ 的 Euler--Poincaré 示性数可以直接用 (1.1.10) 或 (1.1.13) 计算. 其结果虽不甚悦目, 却是

\begin{equation*}
\chi\bigl(\mathbb P^r_k,\Omega^j_{\mathbb P^r_k}(n)\bigr)
=\sum_{0\leq i\leq j}(-1)^{i-j}
\binom{r+1}{i}\binom{r+n-i}{r}.
\end{equation*}

